\documentclass{article}

\usepackage{iclr2027_conference,times}

\usepackage[utf8]{inputenc}
\usepackage[T1]{fontenc}
\usepackage{amsmath}
\usepackage{amssymb}
\usepackage{booktabs}
\usepackage{graphicx}
\usepackage{array}
\usepackage{textcomp}
\usepackage{caption}
\usepackage{subcaption}
\usepackage{placeins}
\usepackage{microtype}
\usepackage{url}
\usepackage[hidelinks]{hyperref}

\newif\iffinalversion
\finalversionfalse
\iffinalversion\iclrfinalcopy\fi

\newcommand{\plotdir}{plots}
% The height cap is a backstop, not a size: it has to stay ABOVE the tallest panel's own
% height or keepaspectratio scales that figure down to meet it and takes its type with it.
% The tallest are the two \halfplot charts at 0.86 aspect -- 0.526\textheight each.
\newcommand{\panelplot}[1]{%
  \includegraphics[width=\linewidth,height=0.55\textheight,keepaspectratio]{#1}}
\let\halfplot\panelplot
\newenvironment{chartfigure}{\begin{figure}[tb]}{\end{figure}}


\title{Zero-Shot Self-Orchestration with Ledger-Based Control \\
Improves Coding in Language Models}

\author{%
  Victor Gao\thanks{Equal contribution.}\ \ \thanks{Persis Capital Inc.} \quad
  Vida Khosrowshahi\footnotemark[1]\ \footnotemark[2] \quad
  Ali Khosrowshahi\footnotemark[1]\ \footnotemark[2] \\
  Xihao Sun\footnotemark[2] \quad
  Juhyun Lee \quad
  Simon (Sang Won) Lee\footnotemark[2] \\
  \texttt{slee@persisholdings.com}
}

\begin{document}

\maketitle

\begin{abstract}
Frontier coding performance is typically bought with larger proprietary models at high
cost. We introduce ledger-based zero-shot self-orchestration, a training-free method in
which fresh instances of one model decompose problems and coordinate through a shared
filesystem holding a plan, notes and current solution. Across nine open and closed-weight
models on the 100 latest \emph{hard} LiveCodeBench problems, the method yields gains of up
to 23.4 percentage points on pinned backends and offers two routes to near-frontier
accuracy. Orchestrated GPT-5.6-Terra reaches 85.0\% pass@1 against Fable 5's 87.4\% at 19\%
of the cost, and locally served, open-weight Qwen3.8-27B rises from 63.0\% to 86.4\%. Gains
are not universal: some models are unchanged or worse. Transcript analysis attributes the
gain to decomposition and persistent context. Inference-time organization can approach
frontier coding accuracy at a fraction of the cost, or on self-hostable weights.
\end{abstract}

\section{Introduction}
\label{sec:intro}

Large language models (LLMs) are increasingly deployed not as single-shot predictors but as
\emph{agents}: models wrapped in a scaffold that lets them plan, call tools, keep notes,
and revise their work over multiple turns. A popular approach to make LLMs is to run \emph{several}
such agents together, on the intuition that a team that decomposes a problem and pools
partial results should beat a single model answering in one pass. That intuition has
produced a large literature \citep{guo2024}, but the evidence is mixed \citep{wangq2024}
and much of it confounded \citep{tran2026}: multi-agent pipelines usually change token
budgets, tool calls, prompts and retrieval at once, so an aggregate gain rarely identifies
which factor helped.

We use \emph{agent} to mean a separately prompted model invocation with a distinct role and
a fresh context; all agents share one underlying model. We call the method studied here
zero-shot self-orchestration: the orchestrator is neither trained for orchestration nor
given task-specific demonstrations of how to decompose or coordinate the problem. Three
properties define it: a shared filesystem workspace --- the \emph{ledger} --- that persists
a plan, notes and current solution across invocations; a manager that re-curates the task
list each round and decides when to stop, rather than following a fixed pipeline; and no
training or per-benchmark tuning, every model role being the same model in a fresh context
with a short generic prompt. Section~\ref{sec:method} describes the loop and
Appendix~\ref{app:scaffold} gives it in full.

On the 100 most recent \emph{hard} LiveCodeBench problems \citep{jain2024}, our training-free
filesystem scaffold improves all three models tested in both arms on a pinned backend by 8.0 to
23.4 points across five paired passes, despite using no router, reward model, search policy,
or benchmark-specific tuning
(Sections~\ref{sec:method} and~\ref{sec:pinned}). Although the scaffold roughly triples the
bill, orchestrating a cheaper model remains the less expensive route to comparable accuracy in
every supported comparison: orchestrated GPT-5.6-Terra trails Claude Fable 5 by 2.4 points
while costing \$11.71 rather than \$61.11 per pass, Qwen3.8-27B comes within 1.0 point, and
GPT-5.6-Luna matches GPT-5.6-Terra's single-call performance at 44\% of the price
(Section~\ref{sec:cost}). These gains diminish as single-call performance rises and can
disappear or reverse: two of the nine models come out flat or negative, with Qwen3.6-35B's
transcripts under reasoning off showing deliberation discarding correct approaches and
unverified worker output overwriting correct work
(Sections~\ref{sec:openrouter} and~\ref{sec:discussion}). We also correct a LiveCodeBench
evaluator bug that penalizes a common input-reading idiom and may therefore affect previously
reported scores (Appendix~\ref{app:evaluator}).

\section{Related work}
\label{sec:related}

Multi-agent systems build on single-agent scaffolds that add structure around one model
\citep{yao2022,shinn2023,madaan2023}; our worker role is an instance of that idea, and what
we vary is the coordination layer above it. Prior systems differ mainly in how that layer is
decided. One line \emph{trains} the coordinator: Fugu \citep{sakana2026} learns to assemble
and direct a pool of expert workers per query, and Conductor \citep{nielsen2025} is a 7B
model RL-trained to design worker communication topologies; earlier query routers learn to
pick a model per query \citep{ong2024}. A second fixes roles and hand-off order in advance,
so every problem flows through the same pipeline
\citep{hong2023,qian2023,li2023,huang2023,islam2024}, while frameworks such as AutoGen
\citep{wu2023} instead supply configurable conversable agents and leave the conversation
pattern open; Trinity \citep{xu2025} sits at the boundary with the first, fixing a
three-role workflow but learning which model fills each role. A third coordinates through a
shared store rather than a pipeline --- ARIADNE \citep{wei2026} pairs a blackboard with MCTS
on program generation, LbMAS \citep{han2025} mediates all communication through a global
board --- or decentralizes entirely, through debate \citep{du2023} or sampling and voting
\citep{li2024}; Mixture-of-Agents \citep{wangj2024a} instead aggregates in layers. A
separate line asks not how coordination is decided but whether reported gains survive
equalizing test-time compute \citep{wangj2024b,tran2026}.

Ours is the training-free counterpart to the first family and the dynamic counterpart to the
second. The manager is a fixed prompt with no learning, which lets us ask how much of the
orchestration benefit is available \emph{without} training, and when it materializes. Its
workspace is a shared store like ARIADNE's blackboard, but coordination is one manager's
plain task-list loop: no MCTS, no reward model, no learned search policy. We do not attempt
an equal-token comparison --- the loop necessarily spends more than a single call --- and
instead ask whether that spend buys better solutions, and whether it buys them more cheaply
than moving to a larger model would, measured throughout against the same model answering in
one call.

\section{Method}
\label{sec:method}

Every role but the verifier is the same underlying model, called in a fresh context and
coordinating through a shared workspace on disk:

\begin{center}
\small
\begin{tabular}{@{}ll@{}}
\texttt{<ws>/task.md}     & the problem statement \\
\texttt{<ws>/plan.md}     & the manager's overarching plan \\
\texttt{<ws>/tasks.json}  & a debug mirror of the task list \\
\texttt{<ws>/notes.md}    & curated ideas / findings / partial proofs \\
\texttt{<ws>/solution.py} & the current solution \\
\end{tabular}
\end{center}

Workers never talk to each other; each sees the ledger plus the one task the manager
assigns. The control flow (Figure~\ref{fig:loop}) is:

\begin{enumerate}
  \item \textbf{Manager --- plan.} The manager reads the problem and writes a 3--6 sentence
  strategy plus 3--6 concrete seed tasks.

  \item \textbf{Worker --- brainstorm (ideation).} The first worker does \emph{not} write a
  solution; it identifies the core difficulty, lists candidate approaches and pitfalls, and
  appends them to \texttt{notes.md}, proposing next steps.

  \item \textbf{Manager --- manage (loop).} The manager folds the plan and the brainstorm
  into one curated task list (merging duplicates, marking done items, adding only genuinely
  new sub-tasks), then either declares the problem done or names the single next task.

  \item \textbf{Worker --- do the task.} A fresh worker executes that one task and rewrites
  both \texttt{solution.py} and \texttt{notes.md} --- the notes are replaced by a curated
  whole rather than appended to --- then proposes remaining steps.

  \item \textbf{Verifier --- run the public tests.} Whenever the round's worker produced a
  fresh candidate, the program is executed against the problem's public stdin tests where it
  has them. The pass/fail verdict, with the first failing case, is fed back to the manager
  and treated as ground truth: a failing run overrides a done verdict, forcing the loop to
  continue with a fix-or-switch task. Control returns to the manager (step 3).

  \item \textbf{Finalizer.} A finalization worker emits the definitive solution whenever the
  loop ends without a clean sign-off --- the round budget is spent, the manager reissues a
  task it just handed out, or it names no task at all. The call is skipped when the manager
  declares the problem done and a usable solution is already on disk, so a redundant final
  pass cannot overwrite a correct answer.
\end{enumerate}

\begin{figure}[tb]
\centering
% The height cap has to stay ABOVE the panel's own height (0.447\textheight here) or
% keepaspectratio scales the figure to fit it and takes the type down with it.
\includegraphics[width=\linewidth,height=0.46\textheight,keepaspectratio]%
  {\plotdir/how_the_manager_arm_answers_one_question_multiagent_py_current_tree_light.pdf}
\caption{\textbf{Manager--worker control flow}, as the v2 scaffold runs it. Boxes labelled
with an agent are model calls in a fresh context; the decisions, the start and end nodes,
and the sample-test run are harness logic and call no model. The boxes absent from the
original scaffold are the verifier of step 5 and the cut-off summarizer.}
\label{fig:loop}
\end{figure}

Guards: a round budget of 10 manager$\to$worker cycles
(\texttt{MAX\_ITERS}); a no-progress guard (if the manager reissues the exact task it just
handed out, the loop stops); and a cut-off summarizer: a worker that hits the token cap
mid-attempt has its partial reasoning summarized by a fresh short call so its ideas still
reach the manager. The plan and notes files are capped when written, at 4,000 and 8,000
characters, so neither grows without bound across rounds.

\paragraph{Baseline.} The single-call baseline is the same model, at the same temperature as
the manager arm's workers (0.2), given the same problem in exactly one call --- no shared
workspace, no loop, and no other role in the prompt. The single call receives the solver
system prompt verbatim; the manager arm's task workers receive that same prompt wrapped in a
subagent preamble and a four-section output contract, and their user message additionally
carries the plan, the current notes and the current solution. The two arms differ only in
that scaffold. Appendix~\ref{app:scaffold} gives the per-role temperatures and the five
differences between this scaffold (v2) and the original one that produced the
OpenRouter-served results of Section~\ref{sec:openrouter}.

\section{Experimental setup}
\label{sec:setup}

\paragraph{Benchmark.} LiveCodeBench code-generation, release\_v6 \citep{jain2024}: the 100
most recent problems in the hard split by contest date, which limits training-data
contamination \citep{jain2024}. We refer to this set as LCB-100. Solutions are graded by
LiveCodeBench's own hidden-test evaluator, after the correction described in
Appendix~\ref{app:evaluator}; every number below is post-correction. The verifier of step 5
runs only stdin-format public tests, which 73 of the 100 problems carry.

\paragraph{Models and serving.} Nine model configurations: Qwen3.5-9B \citep{qwen2026a},
Qwen3.6-35B-A3B \citep{qwen2026b}, Qwen3.8-27B \citep{qwen2026c} (open weights, served
locally in FP8), MiniMax-M3 \citep{minimax2026}, Kimi-K3 \citep{moonshot2026}, and the closed
frontier models Opus-5 \citep{anthropic2026}, the two GPT-5.6 variants Terra and Luna
\citep{openai2026}, and Claude Fable 5 \citep{anthropic2026}. Serving splits the set in two.
The four \emph{pinned-backend} models are served on one fixed backend for v2 comparison: Terra and Luna by the OpenAI API, Qwen3.8-27B by our own vLLM, Fable 5 by the
Anthropic Messages API; because gateway routing proved the dominant source of run-to-run
noise (Section~\ref{sec:discussion}). The five \emph{OpenRouter-served} models ran through a
single OpenAI-compatible gateway on the original scaffold, and are reported separately in
Section~\ref{sec:openrouter} because both the serving path and the scaffold version differ.
Fable 5 was run single-only, so it contributes no manager delta.

\paragraph{Conditions.} The pinned-backend set runs at one setting --- 128k output cap,
reasoning on, $\times 5$ passes --- with Terra, Luna and Qwen3.8-27B in both arms. The cap is
a per-call \texttt{max\_tokens} bound, not a context-window limit. The OpenRouter-served set
runs at three: 128k reasoning-on $\times 1$; 16k reasoning-off $\times 5$; and 128k
reasoning-off $\times 1$ as a control isolating the small cap. Appendix~\ref{app:protocol}
gives the reasoning-effort settings per provider, the instrumentation, and the cap-matching
replay used for Qwen3.8-27B's single arm.

\paragraph{Statistics.} Reported $\pm$SD is the sample SD (divisor $n-1 = 4$) of the five
pass scores. Five-pass comparisons use a paired sign-flip permutation test on per-problem
mean $\Delta$ ($n = 100$ problems; $2\times10^{5}$ resamples), Holm-corrected where a family
of models is tested; 95\% $t$-intervals for mean pass@1 across the five passes
($\mathrm{df} = 4$). Pooled problem$\times$pass agreement uses exact McNemar on the
discordants, treating cells as independent.

\section{Results}
\label{sec:results}

\subsection{Pinned-backend arms, five passes}
\label{sec:pinned}

% Generated by paper_plot_script/make_figures_tex.py -- do not edit.
% The caption text lives in that chart's plot script, next to the numbers it
% describes; re-run the script to update this file.
\begin{chartfigure}
\centering
\halfplot{\plotdir/manager_vs_single_call_four_models_lcb-100_5_passes_128k_max_tokens_reasoning_on_bars_light.pdf}
\caption{\textbf{Manager vs.\ single call.} 128k $\times$ 5 passes, reasoning on. Bars are pass@1 on the same 100 problems, the line through each the 95\% CI across the 5 passes (t, df = 4). The top bar of each pair is the manager, dark; the single call is under it, light; hatch is the model. ``Single call'' is one call with no tools and no loop; Fable 5 ran single-only, so it has one bar, which the dashed rule carries across the chart. The bracket beside each pair is manager $-$ single call; the right-hand column tests the top bar against Fable 5. Both are paired sign-flip permutation tests, unit = problem (n = 100), Holm-corrected within each family of 3: * p $<$ .05, ** p $<$ .01, *** p $<$ .001. The three within-model $\Delta$ all clear p $<$ 1e-4; against Fable 5, p = GPT-5.6-Terra 0.59, GPT-5.6-Luna 4.6$\times$10$^{-4}$, Qwen3.8-27B 0.73. Manager-only vs single-only problem-passes (500 per model): GPT-5.6-Terra 51/11; GPT-5.6-Luna 71/18; Qwen3.8-27B 125/8; exact McNemar 3$\times$10$^{-7}$ or better. The gains are not compensating wins and losses. Every arm is at a 128k cap and re-scored on the corrected evaluator (Appendix C); Qwen3.8-27B's single arm is the 128k cap-matched replay of a 250k generation, which Appendix B describes and Section 5.2 prices.}
\label{fig:main}
\end{chartfigure}


Our headline condition is four models, each run five independent times on LCB-100 at a 128k
cap with reasoning on. Every arm is served on a pinned backend and runs the v2 scaffold except Fable 5. The
earlier runs of the OpenRouter-served set (Section~\ref{sec:openrouter}) cover more models,
but only their 16k reasoning-off condition repeats; the rest are one pass each, through a
gateway whose routing proved the main source of run-to-run noise.

\textbf{All three models with both arms benefit from the manager.} It also narrows
run-to-run spread wherever that spread was wide, for Luna and Qwen3.8-27B alike. Luna's
per-pass $\Delta$ swings across a wide band, so its effect is real but its per-pass estimate
unstable; Terra's flat band is the more reliable. Qwen3.8-27B gains most, its manager arm
reaching Fable 5's single-call range here, and Opus-5's in the OpenRouter-served set. The
direction matches that set's reasoning-on pattern: the smaller the model, the more the
scaffold has to add.

Figure~\ref{fig:main} orders the four models by single-call score: the manager's gain shrinks
monotonically as the single call gets stronger, with two of the manager arms here coming
within a few points of Fable 5's single call, whose lone bar marks where that trend is
heading.

\subsection{What the scaffold costs}
\label{sec:cost}

Every $\Delta$ in the pinned-backend set is bought with tokens. The manager replaces one call
with a plan, a brainstorm, up to ten worker rounds and a finalization pass, and every one of
those is billed. The scaffold roughly triples the bill. Every arm
priced here is at the same 128k output cap, Qwen3.8-27B's single call included
(Appendix~\ref{app:protocol}).

% Generated by paper_plot_script/make_figures_tex.py -- do not edit.
% The caption text and the table numbers live in that chart's plot script;
% re-run the script to update this file.
\begin{chartfigure}
\centering
\halfplot{\plotdir/what_one_pass_costs_lcb-100_5_passes_single_call_vs_manager_against_fable_5_light.pdf}
\caption{\textbf{The scaffold's bill.} Cost of one pass over the same 100 problems; Table~\ref{tab:cost} is the arithmetic behind every bar. No cached-input discount is taken, and Qwen3.8-27B is priced at OpenRouter market rates. The top bar of each pair is the manager, dark; the single call is under it, light; all are at a 128k output cap, and hatch is the model. The right-hand column tests each top bar against Fable 5's single call (Welch per run, $n=5$ vs 5; Holm-corrected across seven comparisons): * $p<.05$, ** $p<.01$, *** $p<.001$. Unbracketed comparisons: each manager $-$ single increase is significant ($8.7\times10^{-5}$ or better per run, $<3.5\times10^{-5}$ per problem); GPT-5.6-Luna's manager undercuts GPT-5.6-Terra's single call by \$1.91 ($2.1\times10^{-7}$ per run, $<3.5\times10^{-5}$ per problem). At 1.18x the assumed Qwen rate (\$0.413/\$3.25 per MTok, inside the spread across hosted providers) the manager's cost advantage over Fable 5 disappears entirely — a larger uncertainty than any p-value here.}
\label{fig:cost}
\end{chartfigure}


\textbf{The scaffold's move is up and to the right in every case.}
Figure~\ref{fig:cost-vs-score} plots price against accuracy for all seven arms, joining the
two arms of each model that has both: every manager arm gains accuracy over its own single
call, and every one costs more. What the extra dollar buys varies by an order of magnitude,
tracking the spread in base prices. The frontier is nearly flat at the top.

\textbf{The base cost varies across models.} A single-call pass on Qwen3.8-27B emits more
than twenty times the output tokens of either OpenAI model: the open model reasons its way to
its score at length; the OpenAI models reach theirs in a fraction of the tokens. Since
reasoning bills as output, a bill here is as much a fact about how much reasoning a provider
lets a model emit as about the scaffold wrapped around it.

\paragraph{Three comparisons along the frontier.}
First, GPT-5.6-Terra with a manager nearly matches Fable 5 for a fifth of the price
(Figure~\ref{fig:cost}). The gap is within error (Figure~\ref{fig:main}), with a one-sided
95\% bound allowing a deficit of 5.8 points, while the saving is significant. Both arms are
priced off API list rates in effect when the experiments were run
\citep{openai2026price,anthropic2026price}.

Second, Qwen3.8-27B with a manager comes within a point of Fable 5's accuracy and costs less.
That difference is again within error, the one-sided 95\% bound allowing a deficit of up to
4.8 points. The cost uses a third-party OpenRouter host rate \citep{openrouter2026}; our test
ran locally, and anyone with capable hardware could host it for less.

Third, the two GPT-5.6 models make the clearest case for spending on the scaffold rather than
on the model: GPT-5.6-Luna with a manager matches GPT-5.6-Terra's single call on accuracy at
under half the price. Luna's lead is not itself significant, but there is no sign of a deficit
either, and the bound rules out its trailing by more than 2.4 points.

\subsection{The effect of truncation}
\label{sec:truncation}

Empty solutions are automatic fails, so any condition that emits fewer of them gains pass@1
without solving more problems. Per-call \texttt{finish\_reason} separates that channel from
genuine problem-solving exactly (Table~\ref{tab:caps}). Neither GPT-5.6 model truncates or
emits an empty answer across all 2{,}000 of their problem-passes, so their $\Delta$s in the
pinned-backend set carry no rescue component. The channel is largest on Qwen3.8-27B, and even
there a cut-off rate is not a loss rate: most truncated generations still contain a complete
solution, written inside the reasoning stream before the cap arrived and recovered from there
by the extractor. Of the single arm's 35 empty cells, the manager solves 25, fails 10 and
leaves none empty, worth 5.0 points, about a fifth of its gain --- a cut-off worker costs a
round rather than the answer. Served locally at a 262{,}144-token context, Qwen3.8-27B's
single arm was in fact generated at 250k and cap-matched to 128k afterwards
(Appendix~\ref{app:protocol}); read at its original 250k cap it gains a little more, by
shrinking the truncation tail rather than lifting the run (Table~\ref{tab:qwen-caps}). That
reading is exploratory, cap-matched to nothing else in the paper, so it cannot carry a
$\Delta$ against the manager arm. Re-scoring each single arm over only the problem-passes
where it emitted code bounds the effect model by model (Table~\ref{tab:emitted}); Fable 5
gains 1.6 points there, six of its nine empty cells being safety refusals rather than
truncations (Section~\ref{sec:discussion}).

\subsection{The OpenRouter-served model set}
\label{sec:openrouter}

The five models below were served through the OpenRouter gateway and run on the original
scaffold, so they are not directly comparable to the pinned-backend arms above.
Table~\ref{tab:openrouter} collects every condition.

\paragraph{Reasoning on.} Opus-5 tops the study, its $+6$ a net of seven problems gained and
one lost. The largest gains come from the smaller models (Figure~\ref{fig:scale-bars}a):
ordered by scale, the single curve rises further than the manager curve, so the scaffold
narrows the spread across model sizes (Figure~\ref{fig:scale-on}, which plots the same numbers
against parameter count).

\paragraph{Reasoning off.} We ran five independent
passes per condition at a 16k cap plus a one-pass 128k control. Run-to-run spread across the
five passes is small. The manager effect is significant for three
of four models, and a Tukey HSD separates Kimi-K3 from all three others and MiniMax-M3 from
Qwen3.6-35B; Qwen3.5-9B sits between those two and separates from neither
(Figure~\ref{fig:scale-bars}b; Figure~\ref{fig:scale-16k} by scale).

The 128k control largely eliminates truncation and reproduces the ordering
(Figure~\ref{fig:scale-bars}c). Empty output is not fully removed --- single-arm emit rates
(non-empty completions) run below the manager's throughout, so part of the gap at the small
end remains an emit-rate effect --- but it cannot explain Kimi-K3, where a 4-point emit gap
accompanies a score gap ten times larger and neither of its arms truncates here at all. Unlike
the reasoning-on setting, the single and manager curves diverge here
(Figure~\ref{fig:scale-off}), though the two settings span different model sets.

\section{Discussion}
\label{sec:discussion}

Table~\ref{tab:transcripts} collects one transcript per model in which the single call failed
and the manager arm solved the problem. Three mechanisms recur. The scaffold commits the
reduction or the bounding argument to disk before any code is written, so the next call builds
on it instead of re-deriving it. It splits coupled objectives into pieces that are implemented
separately. And it bounds each call, so a model that would otherwise talk past its output cap
gets its work written down. In two cases the manager's solution is also markedly smaller than
the failed single-call one.

Truncation is a second place the scaffold earns its keep. State lives in \texttt{plan.md},
\texttt{notes.md} and \texttt{solution.py} rather than in one ever-growing transcript, and
each worker sees a curated view of it --- plan, notes, current solution, one task --- instead
of the full history. The workspace is external memory that the token cap cannot truncate.

That matters because a generation left to run to a very large cap does not merely risk being
cut off; it can degrade the model. Eight of the fourteen Qwen3.8-27B single calls that spent a
full 250k budget without emitting code collapse into repetition --- on \texttt{abc399\_e},
7{,}743 copies of one line --- most of the eight inventing edge cases against a candidate
solution and answering each in turn, hundreds to thousands of times, without terminating. This
is the mechanism behind the otherwise odd finding of Section~\ref{sec:truncation} that most
cut-off generations still contain a complete solution: the model had reached an answer and
then talked past the cap re-checking it. Bounding each call keeps the work far from the cap and
gets it onto disk.

The same argument should extend beyond this benchmark, though we do not test it. A task whose
inputs or intended output do not fit in one call --- a codebase larger than the context window,
a deliverable longer than the output ceiling --- is not solvable in one call at any reasoning
effort, whereas a sequence of bounded steps can carry it.

The manager helps most where the base model is least able to organize its own reasoning, and
that weakness takes two forms. The first is having no internal planning stage at all: with
reasoning disabled a single call jumps straight to code, and the scaffold supplies a planning
space, on disk and across calls, that the model would not otherwise use
(Figure~\ref{fig:scale-bars}c). The second is reasoning that runs away: Qwen3.6-35B with
reasoning on spends its whole budget deliberating, and gains more here than in any of its other
conditions (Figure~\ref{fig:scale-bars}a). Where a model can both plan internally and stop on
its own, the scaffold has less to add.

The manager is not free: its deliberation can settle on a worse answer than a single call would
have produced. On LCB 3765 (Qwen3.6-35B, 128k, reasoning off) the ideation stage identified the
same optimization the single call had used, talked itself out of it as too error-prone, and
committed instead to an asymptotically slower plan that a worker then implemented with a further
bug --- compact, correct code disturbed rather than helped by decomposition. Qwen3.6-35B with
reasoning off is the standout loser on this pattern, negative at both caps.

\textbf{Two inexpensive, training-free additions could address the two halves of that failure},
the plan that was wrong and the implementation that went unchecked. First,
\emph{fresh-perspective workers}. Every worker inherits the notes and the current solution, so a
wrong early approach anchors everything downstream. Spawning some workers with the raw problem
statement and no prior context would give the manager an independent attempt to compare against,
protecting the tight single-call solution the scaffold otherwise disturbs --- the intuition
behind sampling-and-voting and debate \citep{li2024,du2023}. Second, \emph{verification instead
of trust}: the manager currently accepts a worker's report on the basis of a sanity check against
the problem's own samples, so a confidently-wrong solution propagates unchecked and can overwrite
a correct intermediate answer. Running each candidate against a generated test suite before
acceptance would catch regressions of this kind, at the cost of extra calls.

\subsection{Limitations}
\label{sec:limitations}

\begin{itemize}
  \item \textbf{Serving-provider reliability (OpenRouter-served set only).} The five models of
  that set ran through OpenRouter's multi-provider routing, a substantial and
  \emph{time-varying} confound: providers intermittently stalled, dropped the response
  mid-stream, returned error objects, clamped output below the requested cap, or returned
  reasoning-only replies. An attempt with no usable completion is scored a failure, so its
  pass@1 counts infrastructure failures as wrong answers --- a conservative lower bound, and not
  exactly reproducible, since it depends on provider health at run time. Both arms draw on the
  same provider pool, so the single-vs-manager comparison is less sensitive to this than the
  absolute levels are. This confound is what motivated the pinned-backend set, which is why
  those arms and not the OpenRouter sweep carry the paper's headline claims.

  \item \textbf{Provider-served weights may vary (OpenRouter-served set only).} OpenRouter does
  not guarantee one quantization or build per model, adding a second, smaller source of
  run-to-run variation on top of decoding temperature. The headline arms are pinned to one
  backend each, and the locally served Qwen3.8-27B to a specific FP8 checkpoint.

  \item \textbf{Single pass in the OpenRouter-served set.} The reasoning-on and 128k
  reasoning-off comparisons there are one pass per condition, to limit API cost, so their deltas
  are point estimates within the ${\sim}4.5$ percentage points of pass-to-pass variation the 16k
  five-pass runs show on a manager-minus-single delta; per-model significance rests on those
  five-pass runs. It is also why Opus-5's $85 \to 91$, the highest score in the paper, remains a
  point estimate rather than a measured effect.

  \item \textbf{Fable 5 has no manager arm.} It was run single-only, so the strongest
  single-call result in the paper is also the one condition where we cannot say what the scaffold
  would do.

  \item \textbf{Refusals are scored as failures, and only one arm can incur them.} Fable 5
  returned \texttt{stop\_reason=refusal} on 6 of 500 problem-passes, spread over three ordinary
  competitive-programming problems and never deterministic. We deliberately configure no fallback
  model, since a cell labeled Fable 5 must contain Fable 5's outcome, but the consequence is that
  its score carries ${\sim}1.2$ points of classifier penalty no other arm is exposed to, as
  variance rather than a constant offset. Comparisons \emph{against} Fable 5 are therefore mildly
  conservative in its disfavor.

  \item \textbf{Qwen3.8-27B's 128k single arm is a replay, not an independent run.} Those
  generations were produced at 250k and truncated to 128k after the fact
  (Appendix~\ref{app:protocol}). It is the right comparison for isolating the cap --- identical
  generations either side --- but it cannot capture how the model might have budgeted its
  reasoning had it been told the smaller limit up front; we report the as-generated 250k reading
  alongside it (Section~\ref{sec:truncation}).

  \item \textbf{A 9B model with reasoning on was untestable.} Qwen3.5-9B with reasoning enabled
  could not be evaluated in any configuration we tried: through OpenRouter it generates unbounded
  reasoning and returns reasoning-only replies or truncates before an answer appears, and no token
  or effort setting fixed it; locally the checkpoint refused batching, leaving zero completed
  calls. We report it as a model limitation rather than a data point, and it bounds how small a
  ``reasoning'' model this scaffold can be applied to.

  \item \textbf{One benchmark family.} All reported results are competitive-programming code
  generation. Math and knowledge benchmarks (AIME, MATH-500, GPQA, HLE) were run in exploratory
  form but are not reported: on a probe of the ten hardest problems in each, the frontier models
  sit at or near ceiling, leaving little headroom in which to measure a scaffold effect.
\end{itemize}

\section{Conclusion}
\label{sec:conclusion}

With the underlying model held fixed, a training-free manager--worker scaffold over a shared
filesystem improves coding accuracy on every pinned-backend model we ran in both arms. Nothing
in the loop is learned: the coordination is a directory of files that successive invocations of
the same model read and overwrite, driven by one short generic prompt. Transcript analysis
attributes the gain to decomposition and to context that persists across calls.

The scaffold roughly triples the token bill, yet in every comparison our data supports,
orchestrating a cheaper model was the less expensive route to a given accuracy than switching to
a larger one. That leaves two ways to approach frontier coding accuracy without paying frontier
prices: orchestrate a cheaper API model, or orchestrate open weights small enough to serve
locally. The result is conditional rather than general --- the gains shrink as the single call
gets stronger, and some configurations come out flat or worse, where deliberation displaces an
approach that was already right --- but where it holds, inference-time organization substitutes
for model scale.

\subsection*{AI use statement}

In this work, we used generative AI tools as the object of study: every experimental result is a
measurement of language-model output, produced by the scaffold and baselines described in
Section~\ref{sec:method} and released with the code. Separately, generative AI tools assisted
with writing and editing the manuscript, with implementing the orchestration harness and the
analysis and plotting scripts, and with literature search. We did not use generative AI tools to
generate, impute or alter any experimental result, and all reported numbers are computed by the
released scripts directly from the stored run outputs. We have reviewed all AI-assisted work: the
harness and analysis code was read and tested by the authors, the corrected evaluator of
Appendix~\ref{app:evaluator} was found and verified by manual inspection of transcripts, and every
citation was checked against its primary source. We take responsibility for the final content of
this work, including text, claims and artifacts produced with the aid of generative AI.

\subsection*{Ethics statement}

This work involves no human subjects, no personal data and no new dataset release; the evaluation
set is the public \emph{hard} split of LiveCodeBench release\_v6 \citep{jain2024}. Two points bear
on research integrity. First, one arm returned safety refusals that our protocol scores as
failures (Section~\ref{sec:limitations}); we report the rate rather than substituting a fallback
model, so a cell labeled with a model contains that model's outcome. Second, we report a defect in
a widely used public evaluation harness (Appendix~\ref{app:evaluator}) and the corrected
re-scoring of every number, rather than only the numbers that favor our method. The method itself
is a general-purpose inference-time scaffold for code generation and inherits the dual-use profile
of the underlying models; it lowers the cost of capable code generation, which applies to
beneficial and to harmful code alike.

\subsection*{Reproducibility statement}

Section~\ref{sec:method} and Appendix~\ref{app:scaffold} specify the manager--worker loop, its
guards and its per-role temperatures; Section~\ref{sec:setup} and Appendix~\ref{app:protocol} give
the benchmark selection, the serving path and provider settings per model, the output caps and
pass counts, the instrumentation, and the cap-matching replay procedure for Qwen3.8-27B's single
arm. Appendix~\ref{app:evaluator} documents the evaluator defect and the fix, without which the
reported numbers cannot be reproduced from the stored generations. The statistical procedures are
given in Section~\ref{sec:setup}. Two sources of irreducibility are stated in
Section~\ref{sec:limitations}: decoding temperature is non-zero, and the OpenRouter-served arms
depend on provider health at run time. The pinned-backend arms, which carry the paper's claims,
avoid the second.
\iffinalversion
The orchestration scaffold (both versions), the corrected LiveCodeBench evaluator, and the
analysis and figure scripts that produce every number and figure in this paper are available at
\url{https://github.com/slee-persis/GVS5H}, together with all run outputs supporting the results
--- per-problem generations for every arm and condition, per-call \texttt{finish\_reason} and
token logs where instrumented, and the manager arms' workspace transcripts (\texttt{runs/}).
\else
The scaffold, the corrected evaluator, the analysis and figure scripts, and all run outputs
supporting the results --- per-problem generations for every arm and condition, per-call
\texttt{finish\_reason} and token logs where instrumented, and the manager arms' workspace
transcripts --- are provided as anonymized supplementary material, and will be released publicly
on acceptance.
\fi

\iffinalversion
\subsubsection*{Author contributions}

\textbf{[To be completed before the camera-ready.]} For example: V.G., V.K. and A.K.\ conceived
the study and built the scaffold; X.S.\ and J.L.\ ran the experiments and analysed the data;
S.L.\ supervised the work; all authors wrote the paper.

\subsubsection*{Acknowledgments}

\textbf{[To be completed before the camera-ready --- funding sources, if any, and any non-author
help to acknowledge.]}
\fi

\bibliography{references}
\bibliographystyle{iclr2027_conference}

\appendix

\section{The scaffold in full}
\label{app:scaffold}

Section~\ref{sec:method} describes the current scaffold, v2, which produced the pinned-backend
results and is drawn in Figure~\ref{fig:loop}. The original scaffold produced the
OpenRouter-served set of Section~\ref{sec:openrouter}. It runs the same loop on closely matched
prompts, with five differences, all of them on the manager arm: it allows 4 manager$\to$worker
cycles rather than 10 (\texttt{MAX\_ITERS}); it has no public-test verifier at step 5; it has no
cut-off summarizer; its workers append to \texttt{notes.md} rather than rewriting it; and the plan
and notes files feeding a prompt are unbounded rather than capped.

\paragraph{Temperatures.} The manager arm's generative calls use slightly higher temperatures,
where the work is generative rather than executive: 0.3 to write the plan and 0.4 to brainstorm,
against 0.2 for task execution and for curating the task list. The single-call baseline runs at
0.2, matching the workers. These settings are identical in both scaffold versions.

\paragraph{Verifier coverage.} The step-5 verifier runs only stdin-format public tests, which 73
of the 100 problems carry; the other 27 are LeetCode-style, with call-based public tests the
engine does not run, so those get no verifier signal and are checked by the hidden-test grader
alone.

\section{Benchmark and protocol in full}
\label{app:protocol}

\paragraph{Model sizes.} Qwen3.5-9B \citep{qwen2026a} (9B), Qwen3.6-35B-A3B \citep{qwen2026b} (35B
total, 3B active per token), Qwen3.8-27B \citep{qwen2026c} (27B, open weights, served locally in
FP8), MiniMax-M3 \citep{minimax2026} (428B total, ${\sim}23$B active), Kimi-K3
\citep{moonshot2026} (${\sim}2.8$T total, ${\sim}16$ of 896 experts active per token), and the
closed frontier models Opus-5 \citep{anthropic2026}, GPT-5.6-Terra and GPT-5.6-Luna
\citep{openai2026} and Claude Fable 5 \citep{anthropic2026}, all four of undisclosed size.

\paragraph{Caps.} The 128k figure is a per-call \texttt{max\_tokens} bound, not a context-window
limit; for Opus-5, Fable 5 and the GPT-5.6 family it is also the provider's hard ceiling on one
response \citep{anthropic2026,openai2026}. Locally served Qwen3.8-27B has no such ceiling; it is
bounded by its 262{,}144-token context window \citep{qwen2026c} --- confirmed on the FP8 checkpoint
as served (vLLM reports \texttt{max\_model\_len} = 262{,}144) --- shared by prompt and output,
which puts the 250k setting of Section~\ref{sec:truncation} near its largest possible output.

\paragraph{Conditions.} The pinned-backend set runs at one setting: 128k output cap, reasoning on,
$\times 5$ passes, with Terra, Luna and Qwen3.8-27B in both arms and Fable 5 single-only. The
OpenRouter-served set runs at three: 128k reasoning-on $\times 1$ (native
\texttt{reasoning\_effort} where supported, a 20k reasoning budget where not); 16k reasoning-off
$\times 5$, whose repeated passes give run-to-run variance and paired tests; and 128k reasoning-off
$\times 1$, the same without the small cap, isolating it as a confound.

\paragraph{Reasoning.} The output cap and the scaffold are matched across the seven pinned-backend
arms; reasoning depth is not, because no two of these providers expose the same control over it.
Qwen3.8-27B, on our own vLLM, is sent no \texttt{reasoning\_effort} --- the Qwen chat template
leaves reasoning on and unbudgeted, bounded only by the 128k \texttt{max\_tokens}, which on this
stack covers reasoning and answer together. The two GPT-5.6 models are likewise sent none, leaving
them at OpenAI's model default. Fable 5 runs at \texttt{thinking: adaptive} with \texttt{display:
summarized}, and \texttt{output\_config.effort} at \texttt{high}, the API default
\citep{anthropic2026effort}.

\paragraph{Instrumentation.} The two 128k-cap conditions log every call's \texttt{finish\_reason}
and completion-token count, and give each problem a final status in \{\texttt{ok},
\texttt{truncated}, \texttt{empty\_stop}, \texttt{empty}, \texttt{error}\}. The 16k $\times$ 5-pass
runs predate this and record only the extracted code and its pass/fail. An empty code field scores
as a failure.

\paragraph{Cap-matching.} One arm is not natively cap-matched to its partner: Qwen3.8-27B's single
call was generated at 250k while its manager arm ran at 128k (Section~\ref{sec:truncation}). To
compare them at one cap we replay the stored generations rather than re-run the model: each is
truncated to 128,000 output tokens, the solution re-extracted from that prefix, and the result
scored on the same evaluator as every other cell (Appendix~\ref{app:evaluator}). Three details
matter:

\begin{itemize}
  \item \textbf{Truncation is token-exact}, using the serving stack's own tokenizer, not a
  character-count approximation. The deciding cases are those where a complete code block sits just
  before or just after the boundary --- exactly the ones an approximate cut would misclassify.

  \item \textbf{What is truncated is the whole output stream, reasoning then answer}, because that
  is what the original cap bounded. Where the cut lands decides what the extractor sees: past the
  reasoning it sees a truncated answer; inside the reasoning the answer does not exist at all, and
  the harness's empty-content fallback hands it the truncated reasoning instead --- which is why so
  many cut-off generations still yield gradeable code (Section~\ref{sec:truncation}).

  \item \textbf{The cap is invisible to the model, which is what makes the replay exact.} This arm
  ran on our own vLLM, where \texttt{max\_tokens} is enforced by the server: it never enters the
  prompt, and generation simply cuts off at the limit --- 124 of the 500 calls end at
  \texttt{finish\_reason=length}, abruptly rather than wound up. Nor is this local to vLLM: making a
  budget visible to a model takes a separate, opt-in mechanism such as Anthropic's
  \texttt{task\_budget}, while \texttt{max\_tokens} stays an enforced ceiling that truncates the
  response \citep{anthropic2026effort}. The replay therefore does not differ from a genuine re-run
  at 128k.
\end{itemize}

\section{A correction to the LiveCodeBench evaluator}
\label{app:evaluator}

While inspecting transcripts for Section~\ref{sec:discussion} we found a defect in the
LiveCodeBench harness itself, and every number in Section~\ref{sec:results} is reported after
fixing it. The evaluator does not run a candidate as a subprocess; it executes it in-process with
\texttt{sys.stdin} replaced by a mock. That mock's binary view implemented \texttt{readline()} as
\texttt{inputs.split(b"\textbackslash n")[0]} --- a \emph{stateless} expression that returns the
first line on every call. A program reading multi-line input through
\texttt{sys.stdin.buffer.readline()} therefore read line 1 repeatedly and failed however correct it
was. The mock's text-mode \texttt{sys.stdin.readline} was implemented as a proper iterator, and
\texttt{buffer.read()} returns the whole payload, so neither was affected --- which is why the
common \texttt{sys.stdin.buffer.read().split()} idiom never exposed it. We replaced the mock's
binary view with one backed by a \texttt{BytesIO} so that reads advance a position, and re-scored
every stored generation; no model was re-run.

We report it explicitly for two reasons. First, it interacts badly with the v2 verifier: that step
runs the candidate as a \emph{real} subprocess (step 5 of Section~\ref{sec:method}), where
\texttt{buffer.readline()} works correctly, so the manager was told its solution passed the public
samples for programs the grader then marked wrong --- the loop's one external signal certifying a
wrong answer. Second, the exposure is uneven across models rather than a constant offset: 311 of the
3{,}456 pinned-backend outputs containing code use the idiom, and 97\% of those were scored wrong
against 20\% for the outputs that do not. Fable 5 never uses it and its scores are unchanged to the
decimal. The OpenRouter-served set reaches for it once in 5{,}012 generations containing code, and
then through an alias --- Kimi-K3's single call on \texttt{abc396\_g} binds
\texttt{sys.stdin.buffer} to a name and calls \texttt{readline()} on that. Every OpenRouter-served
number is re-scored all the same, and exactly one cell moves: Kimi-K3's 128k reasoning-on single
arm, from 82 to 83. Because the idiom is a matter of coding style, a harness bug of this shape is
not a wash across a leaderboard --- it silently penalizes the models whose style trips it.

\section{Additional results}
\label{app:results}

% Generated by paper_plot_script/make_figures_tex.py -- do not edit.
% The caption text and the table numbers live in that chart's plot script;
% re-run the script to update this file.
\begin{table}[!hb]
\centering
\small
\caption{\textbf{What one pass cost each arm.} List rate $\times$ the tokens it consumed.}
\label{tab:cost}
\begin{tabular}{@{}l@{\hspace{0.7em}}l@{\hspace{0.7em}}r@{\hspace{0.7em}}r@{\hspace{0.7em}}r@{\hspace{0.7em}}r@{}}
\toprule
\textbf{Arm} & \textbf{Rate \$/MTok in / out} & \textbf{In (MTok)} & \textbf{Out (MTok)} & \textbf{\$/pass} & \textbf{\$/solved} \\
\midrule
Qwen3.8-27B single & \$0.35 / \$2.75 & 0.0753 & 7.4247 & \$20.44 & \$0.32 \\
Qwen3.8-27B manager & \$0.35 / \$2.75 & 1.5053 & 18.6277 & \$51.75 & \$0.60 \\
GPT-5.6-Luna single & \$0.20 / \$1.20 & 0.0661 & 0.3299 & \$0.41 & \$0.006 \\
GPT-5.6-Luna manager & \$0.20 / \$1.20 & 1.1686 & 1.0522 & \$1.50 & \$0.019 \\
GPT-5.6-Terra single & \$2 / \$12 & 0.0661 & 0.2728 & \$3.41 & \$0.044 \\
GPT-5.6-Terra manager & \$2 / \$12 & 1.1098 & 0.7911 & \$11.71 & \$0.14 \\
Fable 5 single & \$10 / \$50 & 0.0899 & 1.2043 & \$61.11 & \$0.70 \\
\bottomrule
\end{tabular}
\par\vspace{0.6ex}
{\footnotesize List rates: Qwen3.8-27B \citep{openrouter2026}, GPT-5.6-Luna and GPT-5.6-Terra \citep{openai2026price}, Fable~5 \citep{anthropic2026price}. OpenAI rates are the short-context tier and Anthropic's the base rates; no long-context tier, Batch discount or prompt-caching multiplier is applied.}
\end{table}


% Generated by paper_plot_script/make_figures_tex.py -- do not edit.
% The caption text lives in that chart's plot script, next to the numbers it
% describes; re-run the script to update this file.
\begin{figure}[htbp]
\centering
\panelplot{\plotdir/what_accuracy_costs_lcb-100_5_passes_128k_max_tokens_reasoning_on_all_seven_pinned-backend_arms_light.pdf}
\caption{\textbf{Cost against accuracy.} One point per arm; an arrow runs from the single call to the manager of each model that has both. x is dollars for one pass over the 100 problems (log scale), y is pass@1; the bar through each point is the 95\% CI across the 5 passes (t, df = 4). Light fill = single call, dark = with manager; marker shape is the model. Qwen3.8-27B's single arm is the 128k cap-matched one on BOTH axes - score from the replay, output tokens capped at 128,000 per call to match. Priced as generated at 250k it would sit at \$29.08 rather than \$20.44. Retried and discarded attempts are counted: they were generated and would be billed. The cheapest arm is GPT-5.6-Luna, single call at \$0.41 a pass and the most accurate is Fable 5, single call at \$61.11 - a 149$\times$ spread in price for +20.2 points.}
\label{fig:cost-vs-score}
\end{figure}


\begin{table}[htbp]
\centering
\caption{\textbf{Cap hits and empty solutions per arm.} Single (s) / manager (m), out of 500
problem-passes each. A \emph{cap hit} is a graded generation that reached the token limit
(\texttt{finish\_reason=length}): the single arm's one call, and for the manager its final answer
call only --- internally 108 of its 3{,}235 calls also hit the cap, 95 of them workers, which the
loop absorbs. \emph{No code} is a problem-pass that ended with nothing the extractor could grade.}
\label{tab:caps}
\small
\begin{tabular}{@{}lcccc@{}}
\toprule
\textbf{Model} & \textbf{Cap hits s / m} & \textbf{No code s / m} &
\textbf{of which refusals} & \textbf{Cap} \\
\midrule
GPT-5.6-Terra  & 0 / 0            & 0 / 0          & 0 & 128k\textsuperscript{b} \\
GPT-5.6-Luna   & 0 / 0            & 0 / 0          & 0 & 128k\textsuperscript{b} \\
Qwen3.8-27B    & \textbf{150} / 5 & \textbf{35} / 0 & 0 & 128k\textsuperscript{a} \\
Claude Fable 5 & 3 / --           & 9 / --         & \textbf{6} & 128k\textsuperscript{b} \\
\bottomrule
\end{tabular}

\vspace{0.6ex}
\begin{minipage}{0.92\linewidth}
\footnotesize
\textsuperscript{a}~Every count in this table is at a 128k cap, so the rows are like-for-like.
Qwen3.8-27B's manager arm ran natively at 128k; its single arm was run at 250k and is counted here
cap-matched back to 128k. The 250k counts are given in Table~\ref{tab:qwen-caps}.\\
\textsuperscript{b}~A \emph{hard} ceiling rather than a chosen setting: the GPT-5.6 family and
Fable 5 both cap a synchronous API response at 128k output tokens (Appendix~\ref{app:protocol}), so
for these three arms a truncation cannot be relieved by raising the cap. Qwen3.8-27B, served
locally, has no such ceiling --- which is why its single arm could be run at 250k at all.
\end{minipage}
\end{table}

\begin{table}[htbp]
\centering
\caption{\textbf{Qwen3.8-27B's single arm read at both caps.} Cap-matched to 128k, and as generated
at 250k.}
\label{tab:qwen-caps}
\small
\begin{tabular}{@{}lccc@{}}
\toprule
\textbf{Qwen3.8-27B single} & \textbf{pass@1} & \textbf{Cut off by the cap} &
\textbf{No code at all} \\
\midrule
128k (cap-matched, App.~\ref{app:protocol}) & $63.0 \pm 4.1$ & 150 / 500 \ (30.0\,\%) & 35 / 500 \ (7.0\,\%) \\
250k (as generated)                         & $65.6 \pm 4.6$ & 124 / 500 \ (24.8\,\%) & 21 / 500 \ (4.2\,\%) \\
\bottomrule
\end{tabular}
\end{table}

\begin{table}[htbp]
\centering
\caption{\textbf{Each single arm re-scored} over only the problem-passes where it emitted code.}
\label{tab:emitted}
\small
\begin{tabular}{@{}lccc@{}}
\toprule
\textbf{Model (single arm)} & \textbf{As scored} & \textbf{Emitted code} &
\textbf{Restricted to those} \\
\midrule
GPT-5.6-Terra  & 77.0 & 500/500 & 77.0 \\
GPT-5.6-Luna   & 67.2 & 500/500 & 67.2 \\
Qwen3.8-27B    & 63.0 & 465/500 & 67.7 ($+4.7$) \\
Claude Fable 5 & 87.4 & 491/500 & 89.0 ($+1.6$) \\
\bottomrule
\end{tabular}
\end{table}

\begin{table}[htbp]
\centering
\caption{\textbf{LCB-100 pass@1 (\%), single $\to$ manager.} For the OpenRouter-served models on the
original scaffold. ``--'' = not run; Qwen3.5-9B reasoning-on returns reasoning-only replies and is
unusable. The seven pinned-backend, five-pass arms are reported separately in
Section~\ref{sec:results} and are not pooled here, because both the serving path and the scaffold
version differ (Section~\ref{sec:setup}).}
\label{tab:openrouter}
\small
\begin{tabular}{@{}llccc@{}}
\toprule
\textbf{Model} & \textbf{Params} & \textbf{128k $\cdot$ ON (1 pass)} &
\textbf{128k $\cdot$ OFF (1 pass)} & \textbf{16k $\cdot$ OFF ($\times$5)} \\
\midrule
Opus-5        & n/a            & 85 $\to$ 91 (\textbf{+6})            & --                      & -- \\
Kimi-K3       & ${\sim}2.8$T   & 83 $\to$ 82 ($-1$)                   & 32 $\to$ 74 (\textbf{+42}) & 32.2 $\to$ 62.6 (\textbf{+30.4}) \\
MiniMax-M3    & 428B           & 60 $\to$ 66 (+6)                     & 25 $\to$ 37 (\textbf{+12}) & 21.2 $\to$ 32.2 (\textbf{+11.0}) \\
Qwen3.6-35B   & 35B            & 25 $\to$ 43 (\textbf{+18})           & 35 $\to$ 26 ($-9$)         & 27.8 $\to$ 26.6 ($-1.2$) \\
Qwen3.5-9B    & 9B             & \emph{unusable}                      & 17 $\to$ 20 (+3)           & 14.6 $\to$ 21.8 (\textbf{+7.2}) \\
\bottomrule
\end{tabular}
\end{table}

\begin{table}[p]
\centering
\caption{\textbf{One transcript per model where the single call failed and the manager arm solved
the problem} (Section~\ref{sec:discussion}). LCB problem IDs as in release\_v6; call counts are
manager-arm model calls. Character counts are of the emitted program.}
\label{tab:transcripts}
\footnotesize
\setlength{\tabcolsep}{4pt}
% The four column widths plus the three gaps between them (2\tabcolsep each, so 24pt in
% total) have to fit \linewidth. Stated as fractions rather than fixed cm so that stays
% true: 0.93\linewidth of column and 24pt of gap is 392pt of the 396pt line. The widths
% were 2.3/1.5/4.6/5.4cm, which summed with the gaps to 14.65cm against a 13.97cm line
% and put the table 19.2pt into the right margin.
\begin{tabular}{@{}p{0.155\linewidth}p{0.101\linewidth}p{0.310\linewidth}p{0.364\linewidth}@{}}
\toprule
\textbf{Model} \newline \textbf{(condition)} & \textbf{Problem} & \textbf{How the single call failed} &
\textbf{What the scaffold added} \\
\midrule
Qwen3.5-9B \newline (128k, off) & abc385\_d & Timed out: a na\"ive per-segment scan of a path
simulation. & The brainstorm flagged in writing that the check is $O(N\cdot M) \approx
4\times10^{10}$ and prescribed a range-query structure \emph{before} any code --- the up-front
planning a single pass skips (Section~\ref{sec:discussion}). \\
\addlinespace
Qwen3.6-35B \newline (on) & abc394\_f & Cut off mid-reasoning at 32,768 tokens; returned nothing. &
Solved in one worker cycle: the brainstorm fixed the structural insight (an ``alkane'' subtree needs
a degree-4 center and $\ge 5$ vertices); the worker used an iterative DFS to avoid a recursion-depth
failure. The same 32k clamp was hit 115 times across this model's manager calls, but with several
attempts per problem one clamped call no longer costs the problem. \\
\addlinespace
MiniMax-M3 \newline (on) & 3701 & Printed nothing on \texttt{cdcd}: the lexicographic-reconstruction
branch was broken. & Split the two coupled objectives across worker cycles --- one worker on the
forward cost DP over capped run-lengths, a separate worker on a suffix DP used purely for the
lexicographic reconstruction. \\
\addlinespace
Kimi-K3 \newline (128k, off) & 3687 & Over-counted the longest unique-value path (9 and 3 where the
answers were 6 and 2), never shrinking its window on a repeated value. & The brainstorm wrote the
reduction into \texttt{notes.md} first --- longest substring without repeating characters along each
root-to-leaf path, window start via last-seen depths, $O(n)$ DFS --- and only then did a worker
implement it. \\
\addlinespace
GPT-5.6-Terra \newline (on) & 3688 & 4,615 characters implementing a Segment-Tree-Beats-shaped
structure indexed by candidate value, and wrong. & The plan named the difficulty in advance --- the
state must be compressed rather than maintaining all distinct values per index --- and nine calls
later the manager arm returned 1,860 characters over one standard segment tree. The contribution is
\emph{subtraction}. Manager-only win in four of five passes. \\
\addlinespace
GPT-5.6-Luna \newline (on) & abc397\_e & Carried a \emph{set} of unfinished path lengths per subtree
across 4,092 characters, a general state it never closed. & The brainstorm established the bounding
lemma (at most two unfinished paths at any vertex), after which a worker implemented the postorder
scan directly in 2,166 characters over five calls. Manager-only win in two of five passes. \\
\addlinespace
Qwen3.8-27B \newline (on) & abc397\_g & Spent all 250,000 tokens on edge-case self-interrogation and
never emitted a program --- the runaway-reasoning failure mode of Section~\ref{sec:discussion} in
pure form --- though the reasoning had already reached the min-cut formulation. & Committed that
same reduction to \texttt{plan.md} as a plan, then had a worker implement it: five calls, 3,548
characters, manager-only win in all five passes. The scaffold does not supply the insight; it forces
it onto disk before the budget runs out. \\
\addlinespace
Opus-5 \newline (on) & abc388\_e & The feasibility test over-counted, reporting 225 pairs where 220
is optimal. & Recorded the \emph{justification} in \texttt{notes.md} first --- an exchange argument,
plus monotonicity of feasibility in $K$ licensing a binary search --- then had one worker implement
it and a second harden it. A genuine correction, not truncation relief. \\
\bottomrule
\end{tabular}
\end{table}

% Generated by paper_plot_script/make_figures_tex.py -- do not edit.
% The caption text lives in each condition's plot script (CAPTION/notes) and
% here (the shared lead and the per-model summaries); re-run both to update.
\begin{figure}[p]
\centering
\begin{subfigure}[t]{0.32\linewidth}
  \centering
  \panelplot{\plotdir/manager_vs_single_call_across_model_scale_lcb-100_1_pass_128k_max_tokens_reasoning_on_bars_compact_light.pdf}
  \caption{128k max tokens, reasoning on, one pass per condition.}
\end{subfigure}\hfill
\begin{subfigure}[t]{0.32\linewidth}
  \centering
  \panelplot{\plotdir/manager_vs_single_call_across_model_scale_lcb-100_5_passes_16k_max_tokens_reasoning_off_bars_compact_light.pdf}
  \caption{16k max tokens, reasoning off, five passes per condition.}
\end{subfigure}\hfill
\begin{subfigure}[t]{0.32\linewidth}
  \centering
  \panelplot{\plotdir/manager_vs_single_call_across_model_scale_lcb-100_1_pass_128k_max_tokens_reasoning_off_bars_compact_light.pdf}
  \caption{128k max tokens, reasoning off, one pass per condition.}
\end{subfigure}
\caption{\textbf{Manager vs.\ single call, three conditions.} Models run top to bottom by single-call score, best first, so a model's row differs between panels; the top bar of each pair is the manager. Fill: light = single call (one call, no tools), dark = with manager; $\Delta$ = manager $-$ single, in percentage points. The line through a bar is the 95\% CI across its passes, where the condition ran more than one. A panel prints its two bar values and nothing else; the $\Delta$ and Tukey HSD group letter quoted below per model are what it leaves out (shared letter = not significantly different across the models of that condition, $\alpha = 0.05$). (a) 128k max tokens, reasoning on - effort:high for Kimi-K3 and Opus-5, a 20k reasoning budget requested for Qwen3.6-35B and MiniMax-M3 that providers often did not honor. One pass per condition, so no per-model p. Non-empty completions per 100 (single$\to$manager): Opus-5 97$\to$100, Kimi-K3 97$\to$100, MiniMax-M3 92$\to$93, Qwen3.6-35B 34$\to$68. Empty or truncated output scores as a fail, so $\Delta$ partly tracks emit rate. Opus-5 $\Delta$+6.0 (group ab), Kimi-K3 $\Delta$-1.0 (group b), MiniMax-M3 $\Delta$+6.0 (group ab), Qwen3.6-35B $\Delta$+18.0 (group a). (b) 16k max tokens, reasoning off, five passes per condition; the line through each bar is the 95\% CI across the five passes (t, df = 4). Each $\Delta$ is a paired sign-flip permutation test, unit = problem (n = 100), Holm-corrected across the 4 models; the bracket beside a pair grades its p: * p $<$ .05, ** p $<$ .01, *** p $<$ .001, n.s. otherwise. Kimi-K3 $\Delta$+30.4 (group a), Qwen3.6-35B $\Delta$-1.2 (group c), MiniMax-M3 $\Delta$+11.0 (group b), Qwen3.5-9B $\Delta$+7.2 (group bc). (c) 128k max tokens, reasoning off (ESCALATION\_OR\_REASONING=none), one pass per condition, so no per-model p. Non-empty completions per 100 (single$\to$manager): Qwen3.6-35B 85$\to$95, Kimi-K3 96$\to$100, MiniMax-M3 93$\to$100, Qwen3.5-9B 79$\to$92. Qwen3.6-35B $\Delta$-9.0 (group c), Kimi-K3 $\Delta$+42.0 (group a), MiniMax-M3 $\Delta$+12.0 (group b), Qwen3.5-9B $\Delta$+3.0 (group bc).
}
\label{fig:scale-bars}
\end{figure}


% Generated by paper_plot_script/make_figures_tex.py -- do not edit.
% The caption text lives in that chart's plot script, next to the numbers it
% describes; re-run the script to update this file.
\begin{figure}[htbp]
\centering
\panelplot{\plotdir/manager_vs_single_call_across_model_scale_lcb-100_1_pass_128k_max_tokens_reasoning_on_light.pdf}
\caption{\textbf{Accuracy by model scale, reasoning on, 128k, one pass.} The same numbers as Figure~\ref{fig:scale-bars}a, against parameter count. Dashed = single call, solid = with manager; the block under each tick gives that model's $\Delta$ and Tukey group, and the legend each arm's gain from the smallest model to the largest. }
\label{fig:scale-on}
\end{figure}


% Generated by paper_plot_script/make_figures_tex.py -- do not edit.
% The caption text lives in that chart's plot script, next to the numbers it
% describes; re-run the script to update this file.
\begin{figure}[htbp]
\centering
\panelplot{\plotdir/manager_vs_single_call_across_model_scale_lcb-100_5_passes_16k_max_tokens_reasoning_off_light.pdf}
\caption{\textbf{Accuracy by model scale, reasoning off, 16k, five passes.} The same numbers as Figure~\ref{fig:scale-bars}b, against parameter count; the line through each mark is the 95\% CI across the five passes. Dashed = single call, solid = with manager, and the bracket at a pair grades its Holm-corrected p; the block under each tick gives that model's $\Delta$, p and Tukey group, and the legend each arm's gain from 9B to 2.8T. }
\label{fig:scale-16k}
\end{figure}


\input{fig-128k-reason-off-scale}

\end{document}
